\documentclass[12pt,a4paper]{article}
\usepackage[margin=2.5cm]{geometry}
\usepackage{fontspec}
\usepackage{setspace}
\usepackage{amsmath,amssymb,booktabs}
\usepackage[hidelinks]{hyperref}
\setmainfont{Times New Roman}
\setstretch{1.0}
\setlength{\parindent}{0pt}
\setlength{\parskip}{0.18em}

\begin{document}
\begin{center}
{\Large\bfseries Milestone Report\par}
\vspace{0.4em}
{\large\bfseries Resource-Aware Heuristic Search for Edge Vision Deployment Configuration Selection\par}
\vspace{0.4em}
CS5491 Artificial Intelligence\\
WU JUNYU (72655994, 72655994@cityu-dg.edu.cn)
\end{center}

\textbf{Code repository:} \url{https://github.com/Gene5Wu/rpi-edge-search}

\section{Problem Description and Motivation}

Deploying an object detector on a resource-constrained device involves a
trade-off between detection accuracy, inference latency, and memory consumption.
Testing every configuration is expensive, particularly on edge hardware. This
project asks whether a resource-aware heuristic can identify a high-accuracy
feasible deployment configuration within a limited evaluation budget.

Let \(x=(m,R,p,t)\) denote a configuration consisting of detector variant, input
resolution, precision, and CPU thread count. Given a finite candidate space
\(\mathcal X\), the goal is
\[
\max_{x\in\mathcal X} A(x)
\quad\text{subject to}\quad
L(x)\le L_{\max},\qquad M(x)\le M_{\max},
\]
where \(A\) denotes mAP@0.5, \(L\) median inference latency, and \(M\) peak
process memory. A budget \(B\) limits the number of distinct configurations whose
outcomes a search algorithm may query.

The project is \emph{inspired by} the course topic on Neural Spectral Capacity
(NSC), which uses an inexpensive specification-based score for candidate
selection. Here that idea is adapted to deployment configurations rather than
Transformer architecture design. The claim under test is deliberately narrow: the
score is a heuristic for \emph{ranking and selection}, not a validated predictor
of trained accuracy.

\section{Method Design}

\textbf{Configuration space.}
The space contains YOLOv5n, YOLOv5s, and YOLOv5m, resolutions 320, 416, and 640,
FP32 and INT8 precision, and 1, 2, or 4 CPU threads, giving 54 configurations
(27 when restricted to FP32). Multiple model variants are required: the capacity
term is an architecture quantity, so with a single architecture it would be
constant and could not rank anything. The first measured campaign will prioritize
the 27 FP32 configurations; INT8 remains an extension.

\textbf{Resource-aware heuristic.}
An NSC-inspired capacity term \(C(m)\) is computed from the convolution weight
matrix shapes using a numerical Marchenko--Pastur approximation. The accuracy
proxy is
\[
\widehat A(x)=C(m)\,(R/640)^{\alpha}\,q(p),
\]
and the composite score is
\[
S(x)=\frac{\widehat A(x)}{\max_{z\in\mathcal X}\widehat A(z)}
      -\lambda P_L(x)-\mu P_M(x)-\eta T(x),
\]
where \(P_L\) and \(P_M\) are predicted relative constraint violations clipped at
one, and \(T\) normalizes predicted latency and memory by their candidate-pool
maxima. Latency is estimated from convolution MACs and thread count; memory from
weights, activation size, and runtime overhead. The current assumptions are
\(\alpha=1\), \(q(\mathrm{FP32})=1\), \(q(\mathrm{INT8})=0.97\),
\(\lambda=\mu=1\), and \(\eta=0.05\). These are modelling assumptions, not
empirically fitted constants; the resource model exposes a refit procedure that
will be applied once measurements exist.

\textbf{Search and evaluation.}
Random search samples without replacement. Greedy local search uses observed
outcomes and local neighbours. The proposed best-first method orders its frontier
by \(S(x)\) and never reads the outcome of an unevaluated configuration; the
evaluator enforces this by raising an error for unmeasured rows. Exhaustive
enumeration gives the feasible oracle and is not a budget-matched baseline. The
evaluator counts distinct queries and caches repeats, so the budget reflects
distinct evaluations.

The primary metrics are best feasible accuracy and regret
\(A^*-A_{\mathrm{best}}\). Because two of the algorithms are stochastic, the
constrained comparison is reported over 30 random seeds as mean and standard
deviation, together with the fraction of runs that found at least one feasible
configuration. Rank correlation, the accuracy--latency Pareto front, and
ablations of the capacity and resource terms complete the analysis.

\section{Implementation Progress and Preliminary Results}

The configuration generator, capacity calculation, resource model, scoring
function, budgeted evaluator, four search algorithms, metrics, figure generation
and a multi-seed driver are implemented. A single command,
\texttt{python scripts/reproduce.py}, runs the whole pipeline and currently
reports \textbf{8/8 checks passing}.

\textbf{Benchmark status.} All results below use
\texttt{benchmarks/synthetic\_demo.csv}, a deterministic synthetic table of 54
rows. \emph{These are simulated outcomes, not measured COCO accuracy or Raspberry
Pi performance.} The synthetic truth is generated from partly related analytic
assumptions, so it can favour the heuristic; every figure and table derived from
it is labelled \emph{simulated}. The planned measurement campaign will use a
fixed COCO 2017 validation subset with ONNX Runtime CPU inference.

\begin{center}
\begin{tabular}{@{}llrrrr@{}}
\toprule
Scenario & Budget & Random & Greedy & Heuristic & Oracle \\
\midrule
No resource cap & 4  & 0.4561 & 0.4306 & 0.5503 & 0.5503 \\
No resource cap & 8  & 0.5110 & 0.5110 & 0.5503 & 0.5503 \\
No resource cap & 12 & 0.5206 & 0.5353 & 0.5503 & 0.5503 \\
\bottomrule
\end{tabular}

\small Table 1. Best feasible simulated mAP@0.5, single run, seed 0.
\end{center}

Without resource constraints the heuristic reaches the synthetic oracle with four
queries, while random and greedy need far more.

The constrained picture is different, and it is reported as a distribution rather
than a single run. With a 400 ms latency cap, 22 of the 54 configurations are
feasible and the oracle accuracy is 0.4505.

\begin{center}
\begin{tabular}{@{}lrrrr@{}}
\toprule
Method & Budget 4 & Budget 8 & Budget 12 & Feas.\ success \\
\midrule
Random            & 0.3944 \(\pm\) 0.0394 & 0.4088 \(\pm\) 0.0328 & 0.4171 \(\pm\) 0.0303 & 0.90 / 1.00 / 1.00 \\
Greedy            & 0.3810 \(\pm\) 0.0000 & 0.3952 \(\pm\) 0.0244 & 0.4143 \(\pm\) 0.0307 & 1.00 / 1.00 / 1.00 \\
Heuristic (full)  & 0.4223 & 0.4223 & 0.4223 & 1.00 / 1.00 / 1.00 \\
Heuristic (no capacity term) & 0.3864 & 0.4420 & \textbf{0.4505} & 1.00 / 1.00 / 1.00 \\
\bottomrule
\end{tabular}

\small Table 2. Simulated best feasible mAP@0.5 under a 400 ms cap, 30 seeds for
the stochastic methods (mean \(\pm\) standard deviation) and one deterministic run
for the heuristic. Oracle: 0.4505. Feasible-success rate is given per budget.
\end{center}

Two observations drive the next stage. First, the heuristic is better than or
equal to both baselines at every budget, but its margin over random search
(mainly 0.4223 against 0.4171 at budget 12) is far smaller than in the
unconstrained setting, so the earlier advantage does not transfer automatically
into the constrained regime. Second, and more importantly, \textbf{removing the
capacity term altogether reaches the oracle accuracy} (0.4505) at budget 12.
On this synthetic table the capacity term is therefore not what produces the
constrained improvement.

The ranking evidence points the same way. In the unconstrained pool the full
score has a Spearman correlation of 0.7457 with simulated accuracy. Under the
latency cap its correlation within the feasible subset is 0.6556, against 0.6411
without the resource terms and 0.8534 for predicted MACs alone. A simpler
compute-cost proxy therefore ranks better than the composite score here, which
motivates including a MAC-based baseline in the real comparison rather than
treating the composite score as the strongest available signal.

\section{Reproducibility and Next Steps}

The current pipeline runs on an ordinary laptop:
\begin{verbatim}
pip install -r requirements.txt
python scripts/reproduce.py
python scripts/run_multiseed.py --max-latency 400 --seeds 30
python scripts/reproduce.py --table benchmarks/<measured>.csv --precisions fp32
\end{verbatim}
The scripts generate ranking tables, regret curves, Pareto plots, multi-seed
summaries, and a reproduction summary, without downloading models or requiring a
Raspberry Pi. An FP32-only measured table is a valid benchmark; tools analyse the
covered subset rather than requiring INT8.

The next stage will (1) implement a fixed COCO subset and a real mAP evaluation
with an explicit preprocessing and NMS protocol, (2) collect FP32 latency and
peak-memory measurements on a laptop with isolated per-configuration processes,
(3) repeat the budgeted comparison over multiple seeds including a MAC-based
baseline and the ablations, and (4) examine latency and memory constraints
jointly. Raspberry Pi 4B validation follows the laptop benchmark; INT8 stays an
extension until quantized models can be evaluated.

Two limitations bound these claims. The synthetic benchmark is generated from
assumptions related to the score itself, so it may favour the heuristic and says
nothing about real deployment selection. The analytic architecture specification
also reproduces only 80--84\% of the published parameter counts; the compute
match (93--94\% of published GFLOPs) validates the latency model, but exported
ONNX graphs must be checked before the two are treated as the same network.

\end{document}
